\section{Comparison with Campbell's \texorpdfstring{$q$}{q}-analogue}\label{sec:campbell}

Campbell \cite[Theorem~2]{Campbell1} gives another $q$-analogue of
\eqref{eq:classical}. We show that its summand is an elementary multiple of
the summand of Au's Example~III, which is a modular (not mock modular)
evaluation. For each integer $n\ge0$ put
\[
w_n=1+3q^{2n+1}+3q^{4n+2}-2q^{6n+1}+q^{6n+3}-3q^{8n+2}-3q^{10n+3},
\qquad
u_n=\frac{q^{6n^2}\,w_n\,\poch{q}{q^2}{n}^6}{(1+q^{2n+1})^3\,\poch{q^2}{q^2}{2n}^3}.
\]
The term $u_n$ is the summand of Au's Example~III. In this section only, let
$c_n$ denote Campbell's summand,
\[
c_n=\frac{(-1)^nq^{(n-1)(3n+5)}\poch{q}{q^2}{n}^4\poch{q^3}{q^2}{n}}
{\poch{-q^2}{q^2}{n}^2\poch{q^2}{q^2}{n}^3\poch{q^{3-4n}}{q^2}{n}\poch{q^{2n+2}}{q^2}{n}}
\cdot\frac{\poch{q^{-1},q}{q^4}{n}}{\poch{q^2,q^6}{q^4}{n}}
\cdot\frac{-(1-q^{2n-1})\,w_n}{(1+q^{2n+1})^2}.
\]

\begin{proposition}\label{prop:campbell}
For $0<|q|<1$ and every integer $n\ge0$,
\[
c_n=\frac{1-q^2}{q^6}\,u_n .
\]
\end{proposition}

\begin{proof}
Reversing the order of the factors in the product with negative exponents,
\[
\poch{q^{3-4n}}{q^2}{n}=(-1)^nq^{-3n^2+2n}\poch{q^{2n-1}}{q^2}{n}.
\]
Separating even and odd factors (\cref{fact:evenodd}),
\[
\poch{q^{-1}}{q^4}{n}\poch{q}{q^4}{n}=\poch{q^{-1}}{q^2}{2n},\qquad
\frac{(1-q^{2n-1})\poch{q^{-1}}{q^2}{2n}}{\poch{q^{2n-1}}{q^2}{n}}=(1-q^{-1})\poch{q}{q^2}{n}.
\]
The remaining cancellations use
\[
\poch{-q^2}{q^2}{n}\poch{q^2}{q^2}{n}=\poch{q^4}{q^4}{n},\qquad
\poch{q^2}{q^4}{n}\poch{q^4}{q^4}{n}=\poch{q^2}{q^2}{2n},
\]
\[
\frac{\poch{q^3}{q^2}{n}}{\poch{q}{q^2}{n}}=\frac{1-q^{2n+1}}{1-q},\qquad
\frac{\poch{q^6}{q^4}{n}}{\poch{q^2}{q^4}{n}}=\frac{1-q^{4n+2}}{1-q^2}.
\]
Substituting all of these into $c_n$ leaves $(1-q^2)u_n/q^6$. For $n=0$ all
products are empty and the identity is checked directly.
\end{proof}

Let $F_C$ and $G_C$ be the Wilf--Zeilberger pair used by Campbell
\cite[(41)--(43)]{Campbell1}. He obtains $c_n$ by forming
$F_C(n+1,n+k)+G_C(n,n+k)$ and replacing $q$ by $q^2$. So
\cref{prop:campbell} identifies the transformed summand itself, not only its
classical limit: Campbell's series is, up to the factor $(1-q^2)/q^6$, Au's
modular Example~III, and not the mock modular Watson series studied here.

Campbell's later multiparameter construction \cite[\S4.4, (32)]{Campbell2}
uses, with its base renamed $Q$, the kernel
\[
M_C(n,k)=Q^k\frac{\poch{Q^{n+\alpha},Q^{n+\beta}}{Q}{k}}{\poch{Q^{2n+\gamma},Q^{2n+\delta}}{Q}{k}},
\]
with complex parameters $\alpha,\beta,\gamma,\delta$ (for $0<q<1$ interpret
$Q^\alpha=e^{\alpha\log Q}$). Its ratio $M_C(n,k+1)/M_C(n,k)$ is a rational
function of $Q^k$ whose numerator and denominator have degree at most $2$. For
our kernel \eqref{eq:F}, the ratio $F(n,k+1)/F(n,k)$ is the function
$r^{(k)}$ of \cref{app:cert}; in lowest terms its numerator and denominator
both have degree $6$ in $Q^k$. Multiplying a kernel by a factor depending only
on $n$ does not change this ratio, and specializing the parameters cannot raise
its degree from $2$ to $6$, so the two kernels are not related in this simple way (with the same
base and index). This does not exclude more elaborate transformations
between the two sums.
